\documentclass{beamer}

\date{November 17, 2026}
\author{\L{}ukasz Pi\k{e}tka, Micha\l{} Krzysztosiak, Konrad Miturski, Micha\l{} Sobczak}

\usepackage[T1]{fontenc}
\usepackage[utf8]{inputenc}
\usepackage{lmodern}
\usepackage{graphicx}

\title{Checkpoint 3}

\newcommand{\sequenceframe}[2]{%
  \begin{frame}{#1}
    \centering
    \includegraphics[width=0.98\textwidth,height=0.79\textheight,keepaspectratio]{#2}
  \end{frame}%
}

\begin{document}

\maketitle

\begin{frame}{System communication}
  \begin{itemize}
    \item The customer, staff, and driver interfaces send JSON requests to the Restaurant API under \texttt{/api/v1}.
    \item The backend enforces permissions and coordinates the database, payment provider, and notification and routing services.
    \item Read the sequence diagrams from top to bottom. The labeled alternatives show success, refusal, conflicts, and service failures.
  \end{itemize}
\end{frame}

\sequenceframe{Order and payment}{Sequence Diagrams/restaurant-sequences-04 Order and payment.png}
\sequenceframe{Book a reservation}{Sequence Diagrams/restaurant-sequences-06 Book a reservation.drawio.png}

\begin{frame}{Identifiers and authorization}
  \begin{itemize}
    \item Resource identifiers are UUID strings. Knowing an identifier does not grant access.
    \item Protected routes use \texttt{Authorization: Bearer <token>}. The backend checks the role and, where applicable, ownership or driver assignment.
    \item Guests can submit orders and reservations. The guest access credential for later lookup is still to be specified.
    \item Payment callbacks use a provider signature. The backend verifies the event and its payment reference before changing any state.
  \end{itemize}
\end{frame}

\begin{frame}[fragile]{Order submission: request and contract}
\footnotesize
\begin{verbatim}
POST /api/v1/orders
Idempotency-Key: checkout-4e31
Content-Type: application/json

{
  "items": [{
    "dishId": "123e4567-e89b-12d3-a456-426614174000",
    "quantity": 2
  }],
  "deliveryAddress": {"street": "Example 10", "city": "Warsaw",
                      "postalCode": "00-001", "country": "PL"},
  "contact": {"name": "Example Guest", "phone": "+48123456789"},
  "paymentMethodToken": "pm_example_123"
}
\end{verbatim}
\footnotesize
\texttt{OrderCreate} requires at least one item. Each quantity is an integer from 1 to 100.
The server calculates the total in grosze. A successful submission creates a
\texttt{PENDING\_PAYMENT} order with HTTP 201; it does not confirm payment.
\end{frame}

\begin{frame}[fragile]{Reservation request and booking conflict}
\footnotesize
\begin{verbatim}
POST /api/v1/reservations
Content-Type: application/json

{
  "date": "2026-11-20",
  "time": "18:00",
  "amountOfPeople": 4,
  "contact": {"name": "Example Guest", "phone": "+48123456789"}
}
\end{verbatim}
\footnotesize
\texttt{ReservationCreate} requires a date, local time (\texttt{HH:MM}), a
positive integer party size, and contact information. HTTP 201 returns
the confirmed reservation. If another booking takes the slot first, the
server returns HTTP 409 with the shared response shape
\texttt{error.code}, \texttt{error.message}, and optional \texttt{error.details}.
\end{frame}


\begin{frame}{Decisions for Checkpoint 4}
  \begin{itemize}
    \item Align order state names across the diagrams and API
          (\texttt{PLACED}/\texttt{CONFIRMED}, \texttt{DELIVERED}/\texttt{COMPLETED}).
    \item Specify payment reconciliation after timeouts, late capture after
          cancellation, and refund retry and manual review.
    \item Finalize guest credentials, table duration and timezone, and which
          roles may grant staff permissions.
    \item Check each use case against its sequence and message schema,
          including provider failures and retries.
  \end{itemize}
  \small Full contract: \texttt{api/openapi.yaml}.
\end{frame}

\end{document}
